\newpage 

\section*{About}

\textbf{Introduction}: This text grew out of a reading course at Williams College, MATH 497: Geometry and Quantum Theory, during the Spring 2025 term. The course was based on reading \cite{freed}, and we met once a week for two hours to discuss the material in depth.

These notes were prepared collaboratively by four undergraduate students: Brennan Halcomb, Gary Hu, Rauan Kaldybayev, Theodore Mollano. We gratefully acknowledge the guidance and support of our advisor, Ivo Terek, whose expertise and encouragement greatly enriched this project.

\textbf{Goal}: Daniel Freed discusses quantum and classical physics from a geometric viewpoint. In our notes, we aim to fill in some of the gaps in \cite{freed}, providing additional background, details, and motivation where helpful.

\textbf{Background}: We assume familiarity with functional analysis and differential geometry at the level of introductory graduate courses. A basic understanding of classical and quantum mechanics will also be helpful.

\textbf{Remarks on Specific Chapters}:

\begin{itemize}
    \item \textbf{Chapter 1}: States and Observables

    Freed begins by introducing DOSO (Data of States and Observables), an abstract framework that unifies classical, quantum, and statistical mechanics in a precise, formal way. It specifies some general properties we expect from a physical system. We recommend skimming Lecture 1 initially and returning to it as needed when specific concepts arise in later chapters. This approach lets you focus first on the more concrete and accessible material before revisiting the formal foundations.

    \item \textbf{Chapter 2}: Riemannian and Symplectic Manifolds
    
    This chapter takes a geometric detour to contrast Riemannian and symplectic structures. Rather than developing either theory in depth, we focus on their foundational similarities and differences. The material sets the stage for the symplectic viewpoint on classical mechanics developed in the next chapter.
    
    \item \textbf{Chapter 3}: Hamiltonian and Lagrangian Mechanics
    
    We approach classical mechanics from a geometric perspective, prioritizing geometric structure over physical intuition. Rather than modeling real-world systems, we use symplectic and Riemannian geometry to explore the underlying mathematics. Physical examples appear only to clarify abstract ideas; for a more physics-centered approach, traditional physics texts are recommended.

    \item \textbf{Chapter 4}: Spectral Theorems

    Instead of looking at spectral theorems just from an analytic point of view, we also consider the algebraic and geometric parts to tell a more connected story. We also briefly explore generalizations using harmonic analysis, including ideas like Pontryagin duality.
    
\end{itemize}


%\rauan{Arguably, the first and foundational step in theoretical physics is to translate our intuitive notions such as shape and color into the formal language of mathematics. Greater formalization leads to greater generality and lucidity, albeit at the price of terseness and esotericism. Different physicists and mathematicians make different choices as to how much formalization to pursue. Daniel Freed, perhaps, stands on the more formal side of the spectrum, putting many of the notions left vague by physicists into a precise form and subsuming them under a common formalism. The first step in his project is the definition of DOSO, \emph{Data of States and Observables}, which formalizes the notions of a physical system from classical mechanics, quantum physics, and statistical physics. Roughly speaking, the components of DOSO are:
%\begin{enumerate}
%    \item \textbf{State space.} A \emph{pure state} encodes a possible configuration of a physical system. A \emph{(mixed) state} is a probabilistic object telling us that the system could be in this and that pure state with this and that probability.
%    \item \textbf{Time evolution.} States of the system can be time-evolved in a reasonable way.
%    \item \textbf{Observables.} Certain quantities are accessible to our observation. The space of observables has a nice structure. Each observable specifies a time evolution function.
%\end{enumerate}
%\noindent In quantum mechanics, for example, the state space is the space of density operators; time evolution is specified by the unitary operator $\exp(-i \, t \, H / \hbar)$; and observables are Hermitian operators.}\\

%\noindent \rauan{We axiomatize the state space $\mcS$ as a real convex set (pure states being its extreme points) and a space of potential observables $\mcO$ as a complex topological vector space with structures like a real structure, a Lie algebra on a dense subspace, and a functional calculus. A measurement pairing links observables and states to probability distributions, and a mechanism for dynamics describes transformations generated by observables.}\\

%\noindent \rauan{We illustrate the application of DOSO to classical mechanics and to quantum mechanics. In classical mechanics, pure states correspond to points in phase space, observables to real-valued functions, the Lie algebra to Poisson brackets, and dynamics to Hamiltonian flows. The framework naturally accommodates statistical classical mechanics through probability measures on phase space. In quantum mechanics, states are represented by density operators on a Hilbert space (pure states by rank-one projectors), observables by self-adjoint operators, the Lie algebra by commutators, and dynamics by unitary evolution. The abstract DOSO structure provides a common language compare classical and quantum mechanics.}\\

%\noindent \rauan{We must warn the reader that DOSO is not strictly necessary to understand the remainder of these notes. Indeed, we will largely consider simple special cases of DOSO that can readily be understood without the general definition. Our advice would be to skim Lecture 1 and proceed to the more concrete Lectures 2 and onward. These further lectures explore the mathematics behind some of the most interesting physical ideas and theories. One can then come back to Lecture 1 and appreciate Dan Freed's edifice of an axiomatic formulation of physics.}